% Synthetic test document for h0lon/extract/docs.py (written for the tests, no third-party text).
\documentclass{article}
\usepackage{amsmath,amsthm,amssymb,graphicx}
\newtheorem{thm}{Теорема}[section]
\newtheorem{lemma}[thm]{Лемма}
\theoremstyle{definition}
\newtheorem{defn}{Определение}
\newtheorem*{remark}{Замечание}
\newcommand{\R}{\mathbb{R}}
\newcommand{\Prob}[1]{\mathsf{P}\left(#1\right)}
\title{Случайные величины}
\author{Тестовый автор}
\begin{document}
\maketitle
\begin{abstract}
Учебный пример для проверки извлечения.
\end{abstract}

\section{Определения}

\begin{defn}[Случайная величина]\label{def:rv}
Отображение $X:\Omega\to\R$ называется случайной величиной, если оно измеримо.
\end{defn}

\begin{thm}[Неравенство Чебышёва]
Для любого $\varepsilon>0$ выполнено $\Prob{|X-a|\ge\varepsilon}\le \frac{DX}{\varepsilon^2}$.
\end{thm}

\begin{proof}
Следует из неравенства Маркова.
\end{proof}

\begin{remark}
Замечание без номера.
\end{remark}

Функция распределения
\begin{equation}\label{eq:cdf}
F(x) = \Prob{X\le x},
\end{equation}
где $x\in\R$. Свойства:
\begin{align}
F(-\infty) &= 0,\\
F(+\infty) &= 1 \nonumber
\end{align}

\subsection{Примеры}

\begin{itemize}
\item подбрасывание монеты;
\item бросание кости.
\end{itemize}

\begin{center}
Текст по центру.
\end{center}

\begin{table}
\begin{tabular}{|c|c|}
\hline
$x$ & $P(X=x)$ \\
\hline
0 & 0.5 \\
1 & 0.5 \\
\hline
\end{tabular}
\caption{Закон распределения}
\end{table}

\begin{figure}
\centering
\includegraphics[width=4cm]{img/sample}
\caption{Тестовый рисунок}
\end{figure}

\end{document}
